% ============================================================
% II. THEORETICAL BACKGROUND
% ============================================================

\section{Theoretical background}

\subsection{Sampling and the Nyquist theorem}
\label{sec:aliasing}

An analog signal $x(t)$ is converted into a discrete sequence through periodic sampling \cite{oppenheim}:

\begin{equation}
x[n] = x(nT_s)
\end{equation}

where $T_s = 1/f_s$ is the sampling period. The Nyquist sampling theorem states that a signal can be reconstructed without loss of information only if:

\begin{equation}
f_s \geq 2 f_{\max}
\end{equation}

where $f_{\max}$ is the highest frequency present in the signal. When this condition is not met, \textbf{aliasing} occurs: components with frequency above $f_s/2$ (the Nyquist frequency) ``fold back'' and appear in the reconstructed signal as a component of a different — typically lower — frequency, indistinguishable from a genuine component at that frequency. The standard prevention consists of (a) choosing $f_s \geq 2f_{\max}$ and/or (b) applying an analog anti-aliasing low-pass filter before sampling, attenuating content above $f_s/2$ that the chosen rate cannot represent correctly.

\subsection{Quantization (DAC algorithm)}

Quantization approximates each continuous-amplitude sample to one of $2^{b}$ discrete levels, where $b$ is the converter's bit depth. The \texttt{quantize\_dac} function in \texttt{signal\_tools/sampling.py} implements this process by normalizing the signal to $[0,1]$, rounding to the nearest of $2^b - 1$ intervals, and rescaling back to $[-1,1]$:

\begin{equation}
x_q[n] = \frac{\operatorname{round}\!\left(\dfrac{x[n]+1}{2}(2^b-1)\right)}{2^b-1}\cdot 2 - 1
\end{equation}

A higher number of bits reduces the quantization error but requires a higher-resolution converter.

\subsection{Fast Fourier Transform}

The FFT efficiently computes the Discrete Fourier Transform \cite{lyons}:

\begin{equation}
X[k] = \sum_{n=0}^{N-1} x[n]\,e^{-j2\pi kn/N}
\end{equation}

The \texttt{dominant\_frequency} functions in \texttt{signal\_tools/csv\_signal.py} and \texttt{signal\_tools/audio\_fft.py} compute the single-sided magnitude spectrum via \texttt{numpy.fft.rfft} \cite{numpy}, normalized by $N$, and identify the peak of highest magnitude excluding the $0$ Hz (DC) component.

\subsection{Concurrent sampling mechanisms on a microcontroller}

On a single-core microcontroller (or when all work runs on a single thread), any blocking operation executed within the same loop that performs sampling directly delays the next sample: the \emph{achieved} sampling period stops matching the \emph{nominal} programmed period. The ESP32-WROOM, being a dual-core system with FreeRTOS support and hardware timers, allows two strategies for decoupling sampling from a concurrent compute load:

\begin{itemize}
\item \textbf{Multi-thread/multi-core}: pin a sampling task to a different core than the one running the compute load, so both execute in physical parallel.
\item \textbf{Interrupts}: use a hardware timer that triggers an interrupt service routine (ISR) independently of whatever code is running in \texttt{loop()}, guaranteeing an exact timing reference even if the subsequent processing of each sample is delayed.
\end{itemize}

These two mechanisms are contrasted against the single-threaded blocking baseline in Section~\ref{sec:embebidos}.
